\subsection{向量场在曲线上的积分}

\subparagraph*{Varying force acting on the particle}

\raisebox{-0.5\height}{\includegraphics[height=2cm]{chapters/11/03/0101}}

$\implies$ Total work of $\overrightarrow{F}$ along the path $L=W\approx\sum_{i=1}^n\overrightarrow{F}\left(\overrightarrow{r_i}\right)\cdot\Delta\overrightarrow{r_i}$

$\xlongrightarrow{\text{Take limit as the partition goes to zero}}W\xlongequal{?}\int_L\overrightarrow{F}\,\mathrm{d}\overrightarrow{r}$

\subsubsection[曲线的定向]{Orientation of a Curve}

\raisebox{-0.5\height}{\includegraphics[height=4cm]{chapters/11/03/0102}}

\begin{definition}
    If $L$ can be covered by finitely many local open piece $L_\alpha$ such that $L_\alpha$ admits a regular $C^1$ parametrization $\overrightarrow{r_\alpha}:\,I_\alpha\to L_\alpha$, and on each overlap $L_\alpha\cap L_\beta$, we have $\frac{\mathrm{d}t_\beta}{\mathrm{d}t_\alpha}$ is positive \allbold{(compatible)}, then $L$ is called \emph{orientable}. In this case, the collection $\left\{\left(L_\alpha,\overrightarrow{r_\alpha}\right)\middle|\alpha\in\Lambda\right\}$ of compatible local parametrizations is called an orientation.
\end{definition}

\begin{center}
    \includegraphics[height=5cm]{chapters/11/03/0103}
\end{center}

$\implies$ Unit circle is \emph{orientable}, which is our familiar counterclockwise orientation.

\emph{Equally, an orientation of $L$ is a choice of a continuous unit tangent vector field.}

\raisebox{-0.5\height}{\includegraphics[height=2cm]{chapters/11/03/0104}}\quad Trefoil, Knot Theory

\subsubsection{向量场在曲线上的积分的定义和计算}
\subsubsection{Green 定理}